% FORMALIZACION_NUEVA: prueba explícita de la descomposición y la
% irreducibilidad ya enunciadas en el desarrollo pentacomponente de REV06.
% CERTIFICADO_NUEVO: verificar_caracteres_pentadicos_rev07.py.
% Inserción prevista tras la definición de V12, antes de su descomposición.
\paragraph{Caractères de l’incidence icosaédrique.}
Le calcul suivant s’effectue sur l’action icosaédrique obtenue dans le
théorème~\ref{ii:pent:thm:gravity-exceptional-descent}, avec ses hypothèses
d’incidence et d’orientation. Il détermine la décomposition de la
représentation des sommets et l’irréductibilité de l’espace tensoriel
au moyen de leurs traces, en conservant l’antécédent
\(\pi_{\mathrm{ico}}\) déclaré dans ce théorème.

\begin{proposition}[Décomposition des sommets et caractère tensoriel]
\label{prop:ii-pent-caracteres}
Soient \(\tau=(1+\sqrt5)/2\) et
\(\bar\tau=(1-\sqrt5)/2\). Dans l’ordre des classes
\(1,2A,3A,5A,5B\), les caractères de la représentation rotationnelle
\(V_3\), de sa conjuguée algébrique \(V_{3'}\), de la représentation
des douze sommets \(V_{12}\) et de
\(\operatorname{STF}_3=\operatorname{Sym}^2_0(V_3)\) sont
\begin{equation}
\label{eq:ii-pent-tabla-caracteres}
\begin{array}{c|rrrrr}
\text{classe} &1&2A&3A&5A&5B\\ \hline
\text{cardinal}&1&15&20&12&12\\
\text{classe de }g^2&1&1&3A&5B&5A\\ \hline
\chi_3&3&-1&0&\tau&\bar\tau\\
\chi_{3'}&3&-1&0&\bar\tau&\tau\\
\chi_{12}&12&0&0&2&2\\
\chi_{\mathrm{STF}}&5&1&-1&0&0
\end{array}
\end{equation}
Les représentations \(V_3,V_{3'}\) et
\(\operatorname{STF}_3\) sont absolument irréductibles et distinctes.
De plus,
\begin{equation}
\label{eq:ii-pent-descomposicion-demostrada}
 V_{12}\simeq\mathbf1\oplus V_3\oplus V_{3'}
                  \oplus\operatorname{STF}_3,
 \qquad
 \langle\chi_{12},\chi_{\mathrm{STF}}\rangle=1.
\end{equation}
En particulier, le facteur direct de dimension cinq apparaît une seule fois.
\end{proposition}

\begin{proof}
Dans \(A_5\), les doubles transpositions sont \(5\cdot3=15\) et
les cycles de longueur trois sont \(\binom53\cdot2=20\).
Les \(4!=24\) cycles de longueur cinq se séparent en deux classes de
douze : leur centralisateur est le sous-groupe cyclique d’ordre cinq et est
contenu dans \(A_5\). Si l’on identifie les positions d’un cycle
à \(\mathbb Z/5\mathbb Z\), la permutation qui induit
\(r\mapsto2r\) est un cycle de longueur quatre sur les éléments
non nuls, de signe négatif. Par conséquent, le carré échange les deux
classes. L’inversion \(r\mapsto-r\) est le produit de deux
transpositions et conserve chaque classe. Cela détermine les deux premières
lignes de~\eqref{eq:ii-pent-tabla-caracteres}.

La représentation \(V_3\) est l’action par rotations sur les
coordonnées des sommets. Une rotation d’angle \(\theta\)
a pour trace \(1+2\cos\theta\). Les classes d’ordres deux et trois
donnent respectivement \(-1\) et \(0\). On appelle \(5A\) la
classe d’angles \(\pm2\pi/5\) et \(5B\) celle de
\(\pm4\pi/5\) ; leurs traces sont \(\tau\) et \(\bar\tau\).
Ces expressions s’obtiennent à partir de
\(u^2+u-1=0\), pour \(u=2\cos(2\pi/5)\), et de la formule
de l’angle double. L’équation de \(u\) résulte de la division de
\(1+z+z^2+z^3+z^4=0\) par \(z^2\), avec
\(z=\exp(2\pi i/5)\). On a ainsi calculé le caractère d’une
représentation déjà définie géométriquement.

Pour construire l’autre représentation, on utilise les sommets sans la
normalisation commune \(\nu^{-1}\). Leurs coordonnées appartiennent à
\(K=\mathbb Q(\sqrt5)\). Si \(B\) est la matrice de trois sommets
linéairement indépendants et que \(B_g\) contient leurs images, la
matrice rotationnelle est \(R_g=B_gB^{-1}\in M_3(K)\).
La conjugaison \(\sigma:\sqrt5\mapsto-\sqrt5\) définit alors
\(R'_g=\sigma(R_g)\), et satisfait
\[
 R'_{gh}=R'_gR'_h,\qquad
 (R'_g)^TR'_g=I,\qquad \det R'_g=1.
\]
Par conséquent, \(V_{3'}\) est une représentation réelle effectivement
construite. Son caractère est \(\sigma(\chi_3)\), qui échange
\(\tau\) et \(\bar\tau\).

Le caractère de permutation compte les sommets fixes. L’identité en fixe
douze ; les rotations non triviales d’ordre cinq fixent les deux extrémités
de leur axe. Celles d’ordre deux ou trois ne fixent aucun sommet, car le
stabilisateur d’un sommet est \(C_5\). On obtient ainsi
\(\chi_{12}=(12,0,0,2,2)\).

Pour l’espace tensoriel, si \(\lambda_1,\lambda_2,\lambda_3\)
sont les valeurs propres de \(R_g\), la somme
\(\sum_{i\leq j}\lambda_i\lambda_j\) est
\(\bigl((\sum_i\lambda_i)^2+\sum_i\lambda_i^2\bigr)/2\).
La forme métrique engendre une droite invariante dans
\(\operatorname{Sym}^2(V_3)\), dont le complément sans trace a
pour caractère
\begin{equation}
\label{eq:ii-pent-caracter-stf}
 \chi_{\mathrm{STF}}(g)
 =\frac{\chi_3(g)^2+\chi_3(g^2)}2-1.
\end{equation}
En utilisant l’application de passage au carré du tableau et
\(\tau+\bar\tau=1\), \(\tau\bar\tau=-1\), on obtient
\(\chi_{\mathrm{STF}}=(5,1,-1,0,0)\).

Le produit des caractères réels est
\(\langle\chi,\psi\rangle
=\frac1{60}\sum_C|C|\chi(C)\psi(C)\).
En particulier,
\begin{align*}
 \langle\chi_3,\chi_3\rangle
 &=\langle\chi_{3'},\chi_{3'}\rangle
 =\frac{9+15+12(\tau^2+\bar\tau^2)}{60}=1,\\
 \langle\chi_3,\chi_{3'}\rangle
 &=\frac{9+15+24\tau\bar\tau}{60}=0,\\
 \langle\chi_{\mathrm{STF}},\chi_{\mathrm{STF}}\rangle
 &=\frac{25+15+20}{60}=1.
\end{align*}
La norme un prouve l’irréductibilité des complexifications ; par
conséquent, également l’irréductibilité absolue des représentations
réelles construites. Leurs dimensions et l’orthogonalité distinguent les
trois facteurs directs. Enfin, l’égalité des fonctions centrales
\[
 \chi_{12}=1+\chi_3+\chi_{3'}+\chi_{\mathrm{STF}}
\]
se vérifie dans les cinq colonnes. La semi-simplicité des
représentations d’un groupe fini en caractéristique zéro donne la
décomposition~\eqref{eq:ii-pent-descomposicion-demostrada}. La
multiplicité du facteur direct tensoriel peut se vérifier directement :
\[
 \langle\chi_{12},\chi_{\mathrm{STF}}\rangle
 =\frac{12\cdot5}{60}=1.
\]
\end{proof}

Ainsi, l’idempotent central associé à
\(\chi_{\mathrm{STF}}=\chi_5\) a pour rang
\(5\cdot1=5\) dans \(V_{12}\). L’irréductibilité qui intervient
dans la normalisation de l’entrelaceur tensoriel et la multiplicité qui
intervient dans le couplage gravitationnel sont attestées par
ce même calcul.
